\chapter{EDA IV: Scaling, Transforming Skew, Binning, Encoding and Imbalance}
\companion{scikit-learn User Guide: Preprocessing data}{https://scikit-learn.org/stable/modules/preprocessing.html}

This chapter turns the raw columns that EDA has understood into model-ready features. Each technique comes with the reason it exists, a small numeric
example (computed with scikit-learn/SciPy), and which models need it.

\section{Why scale at all?}
Many algorithms compare features through distances, dot products or gradient steps. If one feature is salary in rupees (range $10^6$) and another is years of
experience (range 30), the salary dominates every distance and creates a stretched loss surface. Scaling puts features on comparable ranges.
\begin{center}\small
\begin{tabularx}{\linewidth}{X X}
\toprule
Needs scaling & Does not need scaling\\
\midrule
$k$-NN, K-Means, SVM (especially RBF), PCA, linear/logistic regression with regularisation, neural networks, anything trained by gradient descent &
decision trees, random forests, gradient-boosted trees (split on order), Naive Bayes (per-feature models)\\
\bottomrule
\end{tabularx}
\end{center}

\section{The scalers, compared on one column}
Column: 10, 20, 30, 40, 1000 (one extreme value).
\begin{center}\small
\begin{tabularx}{\linewidth}{>{\raggedright\arraybackslash}p{2.3cm} l >{\raggedright\arraybackslash}X >{\raggedright\arraybackslash}X}
\toprule
Scaler & Formula & Result & Behaviour\\
\midrule
Standardisation (z-score) & $\frac{x - \mu}{\sigma}$ & $-0.54, -0.51, -0.49, -0.46, 2.00$ & mean 0, sd 1; outlier squashes the rest\\
Min--max normalisation & $\frac{x - \min}{\max - \min}$ & $0, 0.01, 0.02, 0.03, 1$ & range $[0,1]$; very outlier-sensitive\\
Robust scaling & $\frac{x - \text{median}}{\text{IQR}}$ & $-1, -0.5, 0, 0.5, 48.5$ & normal values keep their spread; outlier stays extreme\\
Max-abs & $\frac{x}{\max|x|}$ & $0.01, 0.02, 0.03, 0.04, 1$ & keeps zeros (sparse data friendly)\\
\bottomrule
\end{tabularx}
\end{center}
\subsection{Standardisation vs normalisation: the difference}
\begin{itemize}
\item \textbf{Standardisation} centres to mean 0 and scales to standard deviation 1. No fixed range. Natural for algorithms that assume centred data or use
  variances (PCA, linear models with penalties, SVM, neural networks). Less distorted by outliers than min--max, but still affected.
\item \textbf{Normalisation} usually means \textbf{min--max scaling} to $[0, 1]$. Fixed range; useful when a bounded input is needed (image pixels $/255$, some
  neural-network inputs, distance-based methods on features with known bounds). One outlier compresses all other values near 0.
\item ``Normalisation'' is also used for \textbf{row-wise L2 normalisation}: scale each \emph{row} to unit length, $(3, 4)\to(0.6, 0.8)$, used for TF-IDF vectors
  and embeddings before cosine similarity.
\end{itemize}
\textbf{Robust scaling} (median and IQR) is the choice when outliers exist and you do not want them to dictate the scale.

\begin{trap}
Fit every scaler on the training data only and apply it to test data (Pipeline). Scaling does not change the shape of a distribution: a skewed feature stays
skewed after standardisation; that needs a transformation.
\end{trap}

\section{Curing skew: log, power, Box--Cox, Yeo--Johnson}
\subsection{Why}
Right-skewed variables (salary, price, counts) make linear models chase the few large values, inflate the influence of outliers, produce heteroscedastic
residuals, and make distances dominated by the tail. Transforming toward symmetry often helps linear models, $k$-NN and neural networks (not trees).

\subsection{Log transform}
$x' = \log x$ (or $\log(1 + x)$ when zeros exist). It turns multiplicative differences into additive ones: 1, 2, 4, 8, 16 becomes 0, 0.69, 1.39, 2.08, 2.77: equally
spaced. Skewness $0.89\to0$. Requires $x > 0$ (or $x\ge0$ with \code{log1p}).

\subsection{Other power transforms}
Square root (milder, for counts), reciprocal $1/x$ (strong), square (for left skew).

\subsection{Box--Cox}
A family indexed by $\lambda$, chosen by maximum likelihood to make the data as normal as possible:
\[ x^{(\lambda)} = \begin{cases}\frac{x^\lambda - 1}{\lambda} & \lambda\ne0\\ \log x & \lambda = 0\end{cases}\qquad (x > 0). \]
$\lambda = 1$: no change (shift); $\lambda = 0.5$: square-root-like; $\lambda = 0$: log; $\lambda = -1$: reciprocal-like.
\begin{example}[: Box--Cox]
Fibonacci-like values 1, 2, 3, 5, 8, 13, 21, 34, 55 (skewness 1.25). SciPy picks $\lambda = 0.028$ (almost a log); skewness after $\approx -0.01$.
With $\lambda = 0.5$: 1, 2, 4, 8, 16 become 0, 0.83, 2, 3.66, 6.
\end{example}

\subsection{Yeo--Johnson}
A Box--Cox variant that also accepts zero and negative values (profit/loss, change in salary). scikit-learn's \code{PowerTransformer} uses it by default and
standardises the result. \textbf{Quantile transformer}: maps values to a uniform or normal distribution by ranks; very strong, removes outliers' influence, but
non-linear distortion of distances.

\subsection{Transforming the target}
If you log-transform $y$ for training, predictions must be back-transformed with $\exp$, and $\exp(\E[\log y])$ underestimates $\E[y]$ (Chapter 2): apply a
smearing correction or use a Gamma/Tweedie GLM objective instead. scikit-learn: \code{TransformedTargetRegressor(func=np.log1p, inverse\_func=np.expm1)}.

\section{Numerical to categorical: binning}
\textbf{Why bin?} To capture non-linear effects in linear models (age bands), to reduce the influence of noise and outliers, to make results explainable
(``freshers'', ``mid'', ``senior''), or because the business thinks in bands (salary slabs).
\begin{itemize}
\item \textbf{Equal-width}: split the range into equal intervals (\code{pd.cut}). Simple; sensitive to outliers (most points in one bin).
\item \textbf{Equal-frequency (quantile)}: each bin has the same count (\code{pd.qcut}). Robust to skew; bin edges are data-dependent.
\item \textbf{Domain-based}: experience 0, 1--3, 4--7, 8--12, 13+.
\item \textbf{Clustering-based}: run 1-D K-Means on the values and use cluster IDs as bins (\code{KBinsDiscretizer(strategy="kmeans")}): edges fall in natural gaps.
\item \textbf{Supervised / tree-based}: fit a shallow decision tree of the target on the feature and use its split points (optimal binning, used in credit scoring with
  weight of evidence).
\end{itemize}
\textbf{Costs}: loss of information within bins, arbitrary edges, more parameters. Tree models do not need binning.
\begin{example}[: three binnings of MBA \%]
Values 51.6, 53.3, 55.5, 57.8, 58.8, 59.4, 62.1, 66.3. Edges [0, 55, 60, 100] (domain): counts 2, 4, 2. Quartiles (\code{qcut(q=4)}): 2, 2, 2, 2.
Equal-width into 3 over $[51.6, 66.3]$ (width 4.9): $[51.6, 56.5)$ has 3, $[56.5, 61.4)$ has 3, $[61.4, 66.3]$ has 2.
\end{example}

\section{Categorical to numerical: encodings}
Models need numbers. The encoding must not invent information that is not there (a fake order) and must not leak the target.

\subsection{Label / ordinal encoding}
Map categories to integers. Correct for \textbf{ordinal} variables where order matters: junior $\to$ 0, mid $\to$ 1, senior $\to$ 2. For \textbf{nominal} variables
(city), integers create a fake order and fake distances (``Pune $=$ 2 is between Delhi $=$ 1 and Agra $=$ 3''), which misleads linear models and $k$-NN; tree models
tolerate it somewhat (they can split repeatedly), and it is memory-efficient for very high cardinality in GBMs.

\subsection{One-hot encoding (OHE)}
One binary column per category. City $\in\{$Delhi, Pune, Agra$\}$: Pune $\to(0,1,0)$. No fake order. Drawbacks: many columns for high cardinality (50{,}000 companies),
sparse; unseen categories at test time (\code{handle\_unknown="ignore"}).

\subsection{The dummy variable trap}
With $k$ one-hot columns and an intercept, the columns always sum to 1, which equals the intercept column: \textbf{perfect multicollinearity}. In OLS, $X^\top X$ is singular
and coefficients are not unique (adding $c$ to every city coefficient and subtracting $c$ from the intercept gives identical predictions).
\textbf{Fix}: drop one dummy (\code{drop\_first=True}); the dropped category becomes the \textbf{reference level} and each remaining coefficient is ``the difference from the
reference''. Regularised models (Ridge, Lasso, logistic regression in scikit-learn) and trees do not need the drop (the penalty picks a unique solution), and keeping all
$k$ is more interpretable for them.
\begin{example}[: the trap in numbers]
Rows Delhi, Pune, Agra encoded $(1,0,0)$, $(0,1,0)$, $(0,0,1)$ plus intercept 1: column sum $d + p + a = 1 = $ intercept. Coefficients (intercept 5, Delhi 1, Pune 2, Agra 3)
and (intercept 6, Delhi 0, Pune 1, Agra 2) predict exactly the same.
\end{example}

\subsection{Frequency (count) encoding}
Replace each category by its frequency (or count) in training data. Pune, Pune, Pune, Delhi, Delhi, Agra $\to$ 0.5, 0.5, 0.5, 0.333, 0.333, 0.167. One column whatever the
cardinality; captures ``popularity'' (common companies vs rare); different categories with equal frequency collide.

\subsection{Target (mean) encoding}
Replace each category by the mean of the target for that category. Powerful for high-cardinality features (company, skill, locality), but dangerous:
\begin{itemize}
\item \textbf{Leakage}: computing a row's encoding with its own label lets the model ``see'' the answer; rare categories become perfect predictors. Use \textbf{out-of-fold}
  encoding (each row's value computed from the other folds) or ordered/leave-one-out schemes.
\item \textbf{Noise for rare categories}: smooth toward the global mean:
  \[ \text{enc}(c) = \frac{n_c\,\bar y_c + m\,\bar y}{n_c + m}. \]
\end{itemize}
\begin{example}[: smoothed target encoding]
Placed (1/0) by city: Pune $1,1,0$; Delhi $0,1$; Agra $1$. Global mean $4/6 = 0.667$, smoothing $m = 2$. Agra: raw 1.0 from one student; smoothed $\frac{1\cdot1 + 2\cdot0.667}{3}
= 0.778$. Delhi: raw 0.5, smoothed $0.583$. Pune: raw and smoothed 0.667. Rare categories are pulled toward the global mean.
\end{example}
It is the MAP/empirical-Bayes idea of Chapter 12. scikit-learn's \code{TargetEncoder} cross-fits automatically; CatBoost does ordered target statistics.

\subsection{Other encodings}
\textbf{Binary encoding} (integer code written in binary bits: $\log_2k$ columns), \textbf{hashing} (hash categories into a fixed number of columns; collisions but constant memory;
used in large click models), \textbf{weight of evidence} $\ln\frac{P(\text{cat}\mid y=1)}{P(\text{cat}\mid y=0)}$ (credit scoring), \textbf{learned embeddings} (neural networks: each
category gets a trainable vector; the standard for user/job IDs).

\begin{center}\small
\begin{tabularx}{\linewidth}{l X}
\toprule
Situation & Encoding\\
\midrule
Ordered categories & ordinal\\
Few nominal categories, linear model & one-hot (drop one if unregularised with intercept)\\
Many categories, tree model & native categorical (LightGBM/CatBoost), target or frequency encoding\\
Many categories, linear model & target encoding (out-of-fold, smoothed), hashing, or grouping rare levels into ``Other''\\
Neural network with IDs & embeddings\\
\bottomrule
\end{tabularx}
\end{center}

\section{Imbalanced datasets: the complete toolkit}
Chapter 13 introduced the main ideas; here is the full list the checklist names.
\begin{itemize}
\item \textbf{Random undersampling} (``random deletion''): randomly delete majority-class rows until the ratio is acceptable. Fast; throws away information; good with
  huge data.
\item \textbf{Random oversampling}: duplicate minority rows. Keeps all data; risk of overfitting the duplicated points.
\item \textbf{Weighting the dataset (class weights / sample weights)}: multiply each row's loss by a weight, larger for the minority class
  ($w_c = \frac{n}{K\,n_c}$). No data changed; supported by most models (\code{class\_weight="balanced"}, \code{scale\_pos\_weight}).
\item \textbf{NearMiss} (informed undersampling): keep majority points that are \emph{close} to minority points (version 1: smallest average distance to the 3 nearest minority
  points), so the model focuses on the boundary. Sensitive to noise.
\item \textbf{Tomek links}: pairs of opposite-class points that are each other's nearest neighbours; removing the majority member cleans the boundary. \textbf{Edited nearest
  neighbours} removes majority points misclassified by their neighbours.
\item \textbf{SMOTE}: synthetic minority points by interpolating between a minority point and one of its $k$ minority neighbours: $\mathbf x + u(\mathbf x_{nn} - \mathbf x)$.
\item \textbf{Borderline-SMOTE}: first find minority points ``in danger'': those for which more than half (but not all) of their $m$ nearest neighbours (of any class) are
  majority; generate synthetic points only from those. Points surrounded entirely by majority are treated as noise and skipped. Focuses effort where the boundary is.
\item \textbf{ADASYN}: generate more synthetic points for minority points that have more majority neighbours (harder to learn), in proportion to that difficulty.
\item \textbf{SMOTE-NC} (categorical features), \textbf{SMOTE + Tomek / SMOTE + ENN} (oversample then clean).
\item \textbf{Threshold moving} and \textbf{proper metrics} (PR AUC, recall at precision): often the most effective and simplest step.
\item \textbf{Algorithm-level}: focal loss, balanced random forest, EasyEnsemble (many undersampled subsets, ensembled).
\end{itemize}
\begin{example}[: Borderline-SMOTE danger set]
$m = 5$. Minority point A has neighbours: 3 majority, 2 minority (3/5 $>$ half, not all): \emph{danger}, used for synthesis. Point B: 1 majority, 4 minority: \emph{safe},
skipped. Point C: 5 majority: \emph{noise}, skipped.
\end{example}
Always resample only the training fold (imblearn \code{Pipeline}); never the test set; recalibrate probabilities afterwards if they are used.

\begin{practical}[: where these choices matter in ML at InfoEdge]
Salary engine: log or Box--Cox target, robust scaling of experience, target-encoded company with smoothing. Resdex ranking: frequency-encoded skills, embeddings of titles.
Fraud detection: class weights and threshold tuning rather than heavy SMOTE (fraud patterns change, synthetic points can mislead). Seeker segmentation: robust scaling before
K-Means.
\end{practical}

\stuck{imbalanced-learn User Guide: Over-sampling}{https://imbalanced-learn.org/stable/over_sampling.html}
\section{Interview questions}

\begin{qa}{What is the difference between standardisation and normalisation? When do you use each?}
\oneline{Standardisation rescales to mean 0 and standard deviation 1, with no fixed range; normalisation (min--max) rescales to $[0, 1]$; use standardisation for most
models (linear with penalties, SVM, PCA, neural networks), min--max when a bounded range is needed and there are no extreme outliers, and robust scaling when outliers exist.}
\worked Column 10, 20, 30, 40, 1000: min--max gives 0, 0.01, 0.02, 0.03, 1 (normal values crushed); robust gives $-1, -0.5, 0, 0.5, 48.5$ (normal values keep their spread).
\end{qa}

\begin{qa}{Which models need feature scaling, and why don't tree models?}
\oneline{Distance-, dot-product- and gradient-based models (k-NN, K-Means, SVM, PCA, regularised linear models, neural networks) need it; trees split on thresholds of single features, and
a monotone rescaling does not change which rows fall on each side, so they are unaffected.}
\end{qa}

\begin{qa}{How do you fix a right-skewed feature? Explain Box--Cox.}
\oneline{Apply a log (or $\log(1+x)$), square root, or a fitted power transform; Box--Cox is the family $(x^\lambda - 1)/\lambda$ (log at $\lambda = 0$) with $\lambda$ chosen by maximum
likelihood to make the data as normal as possible; it needs positive data, and Yeo--Johnson extends it to zero and negative values.}
\worked Values 1--55 with skewness 1.25: Box--Cox chooses $\lambda\approx0.03$, skewness $\approx0$.
\end{qa}

\begin{qa}{\deep Does standardising a skewed feature make it normal?}
\oneline{No: standardisation is a linear shift and rescale, so it preserves the shape, including skewness and kurtosis; only non-linear transforms (log, Box--Cox, quantile) change shape.}
\end{qa}

\begin{qa}{What is the dummy variable trap and how do you avoid it?}
\oneline{With one dummy per category plus an intercept, the dummies sum to the intercept column, creating perfect multicollinearity, so OLS coefficients are not unique; drop one category
(the reference level) or omit the intercept; regularised models and trees do not need the drop.}
\end{qa}

\begin{qa}{One-hot vs ordinal vs target vs frequency encoding: when do you use each?}
\oneline{Ordinal for genuinely ordered categories; one-hot for low-cardinality nominal features (especially for linear models); target encoding for high-cardinality features when done
out-of-fold and smoothed; frequency encoding when popularity itself is informative and you need a single compact column.}
\end{qa}

\begin{qa}{\deep Why is naive target encoding dangerous, and how do you do it safely? Compute the smoothed value for a city with one student (placed) when the overall placement rate is
0.667 and $m = 2$.}
\oneline{Because each row's encoding includes its own label, rare categories become perfect fake predictors (leakage and overfitting); compute encodings out-of-fold and smooth toward
the global mean: $(1\cdot1 + 2\cdot0.667)/(1 + 2) = 0.778$.}
\end{qa}

\begin{qa}{How do you convert a numeric feature into categories, and why would you?}
\oneline{Bin it: equal-width, equal-frequency (quantiles), domain-defined bands, clustering-based (1-D K-Means), or tree-based optimal bins; useful to capture non-linearity in linear
models, reduce noise, and produce explainable segments, at the cost of within-bin information.}
\end{qa}

\begin{qa}{List the techniques for imbalanced data and when you would use each.}
\oneline{Class/sample weighting (first choice, no data change), random undersampling or ``random deletion'' (very large data), random oversampling, SMOTE and its variants
(Borderline-SMOTE near the boundary, ADASYN for hard regions), NearMiss and Tomek links (informed undersampling/cleaning), ensembles of undersampled sets, threshold tuning and
the right metrics.}
\end{qa}

\begin{qa}{How does Borderline-SMOTE differ from SMOTE and ADASYN?}
\oneline{SMOTE interpolates from all minority points; Borderline-SMOTE only from minority points ``in danger'' (more than half but not all neighbours are majority), ignoring safe and noisy
points; ADASYN creates more points for minority samples with more majority neighbours, in proportion to their difficulty.}
\end{qa}

\begin{qa}{\deep After class weighting your model's predicted probabilities are too high. Why?}
\oneline{Weighting (like oversampling) changes the effective class prior seen in training, so the model learns an inflated base rate; the ranking is fine but probabilities must be recalibrated
on unweighted validation data (or the intercept corrected) if used as probabilities.}
\end{qa}

\begin{qa}{What is row-wise L2 normalisation and where is it used?}
\oneline{Dividing each row vector by its length so it has norm 1, e.g. $(3,4)\to(0.6,0.8)$; used for TF-IDF and embedding vectors so that dot products equal cosine similarities.}
\end{qa}
